\scanpage{251}
\mspage{125}
% archive sheet number: 250

flèche $f:X\unsure{-|\to}X'$ dans $\mathrm{Cat}_{/Y}$, on a $f\in W$ ssi $\forall i\in I$, \deleted{\unsure{on a}} le morphisme $\unsure{f^i|f_i}:X_i-X'_i$ au dessus de $Y_i$ est dans $W$.

\uline{Corollaire 9.6}. Si l'axiome Loc $(\deleted{\unsure{6 bis}}5)$ est satisfait, alors l'axiome de connexité Loc $(8)$ équivaut à ceci : pour toute flèche $f:X-Y$ de $W$, avec $Y$ \unsure{$0$-connexe|$o$-connexe}, $X$ est $0$-connexe.

NB Les axiomes Loc $(8)$ (connexité) et Loc $(\deleted{\unsure{6}}5\,\mathrm{bis})$ \deleted{\unsure{ou Loc (8 bis)}} (de décomposition) ont l'air de nature voisine, \deleted{\unsure{je dirais $=$ les avantages,}} \deleted{\unsure{car}} étant liés l'un et l'autre \deleted{\unsure{ou que}} à la décomposition des objets de $\mathrm{Cat}$ en composantes connexes. (Néanmoins, je pense qu'ils sont indépendants.) Je présume qu'on n'aura jamais à travailler vraiment avec des $W$ qui ne satisfont à ces deux axiomes ! (En plus de Loc $(1,2,3)$, qui en constituent le ``minimum vital''.)

Mais apparemment il n'y en a pas besoin dans les développements théoriques généraux jusqu'ici. (Mais l'axiome de connexité Loc $(8)$ pourrait jouer un rôle, quand on parle de la \unsure{déterminabilité|détermination} réciproque des structures de \unsure{contractibilité|contractilité}, et de $W$-asphéricité, cf \unsure{P.S.})
